\section{Прямая сумма подпространств. Теорема о разложении ЛП в прямую сумму подпространств}

\begin{definition}
Сумма подпространств $U_1, U_2, \dots, U_k$ называется \textbf{прямой} (обозначается $U_1 \oplus U_2 \oplus \dots \oplus U_k$), если любой вектор $x$ из этой суммы единственным образом представляется в виде суммы:
\[
x = x_1 + x_2 + \dots + x_k, \quad \text{где } x_i \in U_i \text{ для } i = 1, \dots, k.
\]
\end{definition}

\begin{theorem}[О разложении пространства в прямую сумму]
Пусть конечномерное линейное пространство $V$ разлагается в прямую сумму подпространств: $V = U_1 \oplus U_2 \oplus \dots \oplus U_k$. Если в каждом подпространстве $U_i$ выбрать по базису, то объединение этих базисов образует базис всего пространства $V$.
\end{theorem}

\begin{proof}
Пусть для каждого $i = 1, \dots, k$ система $E_i = \{e_{i1}, e_{i2}, \dots, e_{im_i}\}$ является базисом подпространства $U_i$. Рассмотрим объединенную систему векторов $E = E_1 \cup E_2 \cup \dots \cup E_k$. Докажем, что $E$ --- базис $V$, проверив два условия: порождение и линейную независимость.

\begin{enumerate}
    \item \textbf{Система $E$ порождает $V$:} 
    Возьмем произвольный вектор $x \in V$. Так как $V = U_1 \oplus \dots \oplus U_k$, вектор $x$ единственным образом представляется в виде $x = x_1 + x_2 + \dots + x_k$, где $x_i \in U_i$. 
    Поскольку $E_i$ --- базис $U_i$, каждый $x_i$ линейно выражается через векторы $E_i$:
    \[
    x_i = \sum_{j=1}^{m_i} \lambda_{ij} e_{ij}.
    \]
    Следовательно, $x = \sum_{i=1}^k \sum_{j=1}^{m_i} \lambda_{ij} e_{ij}$, то есть $x$ линейно выражается через объединенную систему $E$.

    \item \textbf{Система $E$ линейно независима:} 
    Составим линейную комбинацию всех векторов системы $E$, равную нулевому вектору:
    \[
    \sum_{i=1}^k \left( \sum_{j=1}^{m_i} \lambda_{ij} e_{ij} \right) = 0.
    \]
    Обозначим $v_i = \sum_{j=1}^{m_i} \lambda_{ij} e_{ij}$. По построению, $v_i \in U_i$. Тогда мы получили разложение нулевого вектора:
    \[
    v_1 + v_2 + \dots + v_k = 0.
    \]
    С другой стороны, нулевой вектор тривиально раскладывается как $0 + 0 + \dots + 0$. Поскольку сумма $U_1 \oplus \dots \oplus U_k$ прямая, разложение любого вектора (в том числе нуля) единственно. Следовательно, $v_i = 0$ для всех $i = 1, \dots, k$.
    
    Но $v_i$ --- это линейная комбинация базисных векторов подпространства $U_i$. В силу линейной независимости базиса $E_i$, все коэффициенты $\lambda_{ij}$ должны быть равны нулю. 
    
    Таким образом, из равенства нулю линейной комбинации векторов системы $E$ следует, что все ее коэффициенты нулевые, что и означает линейную независимость системы $E$.
\end{enumerate}
Система $E$ линейно независима и порождает $V$, следовательно, она является базисом $V$.
\hfill$\blacksquare$
\end{proof}

\begin{corollary}
Если $V = U_1 \oplus U_2 \oplus \dots \oplus U_k$, то размерность пространства $V$ равна сумме размерностей подпространств:
\[
\dim V = \dim U_1 + \dim U_2 + \dots + \dim U_k.
\]
\end{corollary}
